
        \documentclass[fleqn]{article}      
        \usepackage{amsmath}
        \usepackage{fontspec}
        \usepackage{kotex} % 한국어 지원  

        \begin{document}      
         객관식 질문

1. 다음 중 적분의 정의는 무엇인가요?
   a) 미분의 역연산  
   b) 함수의 기울기를 구하는 것  
   c) 함수의 면적을 구하는 것  
   d) 방정식을 푸는 것  

2. 다음 중 정적분의 기호는 무엇인가요?
   a) ∂  
   b) ∫  
   c) ∑  
   d) ∏  

3. ∫(x^2)dx의 결과로 나오는 항은?
   a) (1/3)x^3 + C  
   b) (1/2)x^2 + C  
   c) (1/4)x^4 + C  
   d) (1/5)x^5 + C  

4. ∫(sin x)dx의 결과는 무엇인가요?
   a) -cos x + C  
   b) sin x + C  
   c) -sin x + C  
   d) cos x + C  

5. 다음 중 정적분의 기본 정리에 대한 설명으로 옳은 것은?
   a) 미분과 적분은 서로 연관이 없다.  
   b) 정적분은 함수의 기울기를 나타낸다.  
   c) 미분과 적분은 서로 역의 관계에 있다.  
   d) 정적분은 함수의 최대값을 찾는 방법이다.  

 주관식 질문

1. 적분의 정의를 설명하시오.
2. 정적분과 부정적분의 차이를 설명하시오.
3. ∫(3x^2)dx를 계산하시오.
4. ∫(e^x)dx의 결과를 구하시오.
5. 기하학적 의미에서 적분이 무엇을 나타내는지 설명하시오.
6. 기본 정적분의 정리를 설명하시오.
7. ∫(1/x)dx의 결과는 무엇인가요?
8. 함수 f(x) = x^3에 대한 정적분을 [1,2] 구간에서 계산하시오.
9. 정적분의 의미를 그래프를 통해 설명하시오.
10. ∫(tan x)dx의 결과를 구하시오.
11. 적분의 수치적 근사 방법 중 하나를 설명하시오.
12. 함수의 적분이 연속 함수임을 증명하시오.
13. 적분의 변수가 바뀌었을 때의 처리 방법을 설명하시오.
14. 특정한 구간에서의 정적분과 그 구간의 기하학적 의미를 설명하시오.
15. 수치적 적분을 위한 사다리꼴 법칙에 대해 설명하시오.
16. 적분을 이용한 면적 계산의 예를 제시하시오.
17. ∫(x^2 + 3x + 2)dx를 계산하시오.
18. ∫(cos x)dx의 결과를 구하시오.
19. 적분을 사용하여 물리적 현상을 설명하시오.
20. ∫(x^n)dx의 일반적인 결과를 구하시오.
21. 적분의 연속성에 대해 설명하시오.
22. 중첩 정적분의 정의를 설명하시오.
23. 다변수 함수의 적분을 설명하시오.
24. 함수의 면적을 구하기 위한 적분의 기초 개념을 설명하시오.
25. 적분의 응용 분야를 나열하시오.

 쉬운 문제

1. ∫(x)dx를 구하시오.
2. ∫(1)dx의 결과는 무엇인가요?
3. ∫(x^3)dx를 계산하시오.
4. ∫(2x)dx의 결과를 구하시오.
5. ∫(5)dx는 얼마인가요?
6. ∫(sin x)dx의 결과를 구하시오.
7. 정적분의 기호는 무엇인가요?
8. ∫(cos x)dx의 결과는 무엇인가요?
9. ∫(x^2 + 1)dx를 계산하시오.
10. ∫(3x^2 + 2x)dx의 결과를 구하시오.

 중간 문제

1. ∫(x^4)dx를 계산하시오.
2. ∫(e^x)dx의 결과를 구하시오.
3. ∫(1/x)dx의 결과는 무엇인가요?
4. ∫(tan x)dx의 결과를 구하시오.
5. ∫(x^3 - 2x + 1)dx를 계산하시오.
6. 특정 구간 [0,1]에서 ∫(x^2)dx의 결과를 구하시오.
7. ∫(sin^2 x)dx의 결과는?
8. ∫(x^2 * e^x)dx를 계산하시오.
9. ∫(x^2 + 3)dx의 결과를 구하시오.
10. ∫(x^3 + 2x^2 - x + 5)dx를 구하시오.

 어려운 문제

1. ∫(x * ln x)dx를 계산하시오.
2. ∫(x^2 * sin x)dx의 결과를 구하시오.
3. ∫(e^(2x) * cos(3e^x))dx를 계산하시오.
4. ∫(x^2 * ln(x))dx를 구하시오.
5. ∫(1/(1+x^2))dx를 계산하시오.
6. ∫(x^3 * e^(-x^2))dx의 결과를 구하시오.
7. ∫(cos^2 x)dx의 결과를 구하시오.
8. ∫(x^n * e^x)dx를 구하시오.
9. ∫(tan^2 x)dx의 결과는 무엇인가요?
10. ∫(x^3 * cos x)dx를 계산하시오.

      
        \end{document} 
    